\documentclass[../Main.tex]{subfiles}

\begin{document}
\section{Spanning Sets}
\begin{definition}{Span}
    Let $V$ be a vector space. The \underline{span} of a subset $S \subseteq V$, $\spn{S}$, is the set of elements that can be written as:
    \begin{equation*}
        \finitesum_{s \in S} \lambda_s s
    \end{equation*}
\end{definition}
Here we use finite sum notation: all but finitely many $\lambda_s$ can be non-zero. This sum is called a \underline{linear combination} of elements of $S$.
\begin{definition}{Spanning set}
    A subset $S$ of a vector space $V$ is a \underline{spanning set} if $\spn{S} = V$.
\end{definition}
\begin{propositions}{
        
        \label{propsSpanning}
    }
    \item For $S \subseteq V$, $\spn{S} \leq V$. Further, if $S \subseteq W \leq V$ then $\spn{S} \leq W$.
    \item If $S, T \subseteq V$ and $S$ spans $V$, if $S \subseteq \spn{T}$ then $T$ also spans $V$.
    \item $\spn{S \cup T} = \spn{S} + \spn{T}$.
\end{propositions}
We also have that by convention, $\spn{\emptyset} = \{\zv\}$.
\section{Linear Independence}
\begin{definition}{Linear independence}
    A subset $S$ of a vector space $V$ is \underline{linearly independent} if:
    \begin{equation*}
        \finitesum_{s \in S} \lambda_s s = \zv \implies \lambda_s = 0~\forall s\in S
    \end{equation*}
\end{definition}
Note that every subset of a linearly independent set is linearly independent.
\section{Basis of a Vector Space}
\begin{definition}{Basis}
    A subset $S$ of a vector space $V$ is a  \underline{basis} for $V$ if $S$ is linearly independent and spans $V$.
\end{definition}
\begin{proposition}[Finite-dimensionality in terms of bases]
    Let $S \subseteq V$ be a finite spanning set for $V$. Then there exists $S' \in V$ such that $S'$ is a finite basis for $V$.
    \label{propFiniteSpanFiniteBasis}
\end{proposition}
\begin{proof}
    Initially set $S' = S$. Repeatedly remove vectors that have a zero linear combination until $S'$ is linearly independent. The process must terminate because $S$ is finite.
\end{proof}
\begin{corollary}
    Every finite-dimensional vector space has a finite basis.
    \label{corFinDimFinBasis}
\end{corollary}
\begin{theorem}[Steinitz Exchange Lemma]
    Let $S, T$ be finite subsets of a vector space $V$. Let $S$ span $V$ and let $T$ be linearly independent. Then:
    \begin{enumerate}
        \item $|S| \leq |T|$;
        \item There exists $T' \subseteq T$ such that $|T'| = |T| - |S|$ and $S \cup T'$ spans $V$.
    \end{enumerate}
    \label{thmSteinitzExchange}
\end{theorem}
\begin{proof}
    Label the elements of $T$ and $S$. If $S$ is empty then done so assume not. Write the first vector of $S$ as a linear combination of elements of $T$. Choose the first element of $T$ with a non-zero coefficient in this linear combination, and make it the subject. Then we can remove this vector and replace it with the first element of $S$ such that the new set $S_1$ still spans $V$.

    At step $k$, do the same but now we need to make sure that there is an element of $T$ in $S_k$ that has a non-zero coefficient. Do this by discounting the case that all the elements of $T$ have zero coefficients in the linear combination by linear independence of $S$.

    Conclude at step $|S|$ that we have the required results. We show $|S| \leq |T|$ by contradiction (a spanning set cannot live inside a linearly independent set).
\end{proof}
\begin{corollary}
    Let $V$ be finite dimensional. Then every basis is finite and has the same size.
    \label{corBasisSameSize}
\end{corollary}
\begin{proof}
    Let $B$ be a basis. Prove by contradiction using \thmref{thmSteinitzExchange} that another basis $B'$ must be finite. Now apply \thmref{thmSteinitzExchange} with $S = B$ and $T = B'$, then the other way around, to show that $|B| = |B'|$.
\end{proof}
\begin{definition}{Dimension}
    The \underline{dimension} of a vector space $V$, $\dim(V)$, is the size of any basis of $V$.
\end{definition}
\begin{corollary}
    Let $V$ be a finite-dimensional vector space and $U \leq V$. Then $\dim(U)\leq\dim(V)$ with equality if and only if $U = V$.
    \label{corSubspaceDim}
\end{corollary}
\begin{proof}
    First find a basis of $U$ by induction. Then use \thmref{thmSteinitzExchange} for the rest.
\end{proof}
\begin{proposition}[Basis Extension Principle]
    Let $V$ be a finite-dimensional vector space and $U \leq V$. For any basis $B_U$ of $U$, there exists a basis $B_V$ of $V$ such that $B_U \subseteq B_V$.
    \label{propBasisExtend}
\end{proposition}
\begin{proof}
    Let $T$ be any basis for $V$ and apply \thmref{thmSteinitzExchange} with $S = B_U$ to get $B_V = B_U \cup T'$. Prove by considering sizes that this is indeed a basis.
\end{proof}
\section{Bases and Linear Maps}
In this section, let $\alpha : V \mapsto W$ be a linear map, $V, W$ vector spaces over $\F$.
\begin{definition}{Nullity}
    The \underline{nullity} of a linear map, $n(\alpha)$ is the dimension of the kernel.
\end{definition}
\begin{definition}{Rank}
    The \underline{rank} of a linear map, $\rk(\alpha)$ is the dimension of the image.
\end{definition}
\begin{lemma}
    Let $V$ be finite-dimensional and $U \leq V$. Then:
    \begin{equation*}
        \dim(V / U) = \dim(V) - \dim(U)
    \end{equation*}
    \label{lemDimQuotient}
\end{lemma}
\begin{proof}
    Let $B_U$ be a basis for $U$ and use \thmref{propBasisExtend} to find a basis $B_V$. Label the elements of each. Create the set:
    \begin{equation*}
        B_{V / U} = \{v_i + U\}
    \end{equation*}
    where $v_i$ are the basis vectors of $V$ but not $U$. Show this is a basis of $V / U$.
\end{proof}
\begin{theorem}[Rank-Nullity Theorem]
    Let $V$ be finite-dimensional. Then:
    \begin{equation}
        \dim(V) = n(\alpha) + \rk(\alpha)
        \label{eqnRankNullity}
    \end{equation}
    \label{thmRankNullity}
\end{theorem}
\begin{proof}
    Combine \thmref{thmIsomorphism} and \thmref{lemDimQuotient} for the result.
\end{proof}
\begin{corollary}[Linear pigeon-hole principle]
    Assume that $\dim(V) = \dim(W) = n$. The following are equivalent:
    \begin{enumerate}
        \item $\alpha$ is injective;
        \item $\alpha$ is surjective;
        \item $\alpha$ is an isomorphism.
    \end{enumerate}
    \label{corLinPigeon}
\end{corollary}
\begin{proof}
    By definition statement 3 implies 1 and 2. Prove that 1 implies 2 and 2 implies 1 by \thmref{thmRankNullity}.
\end{proof}
\begin{proposition}
    Let $V$ be a (perhaps infinite-dimensional) vector space with basis $B_V$. For any function $f : B_V \mapsto W$ where $W$ is any vector space, there exists a corresponding unique linear map:
    \begin{equation*}
        F : V \mapsto W
    \end{equation*}
    such that $F(b) = f(b)$ for all basis vectors $b$.
    \label{propLinMapExtend}
\end{proposition}
\begin{proof}
    Define $F$ by applying $f$ to linear combinations. For a basis, linear combinations are unique so $F$ is well-defined (can be shown by contradiction). Now show linearity, and uniqueness by again considering $F$ and $G$ applied to linear combinations.
\end{proof}
\begin{corollary}
    Let $\dim(V) = n$ and $B$ a basis for $V$ with elements $\{b_i\}_{i = 1}^n$. Then there exists an isomorphism:
    \begin{align*}
        F_B : V &\mapsto \F^n\\
        \sum_{i=1}^n \lambda_i b_i &\mapsto
        \begin{pmatrix}\lambda_1 \\ \vdots \\ \lambda_n\end{pmatrix}
    \end{align*}
    \label{corCoordinateVectorMap}
\end{corollary}
\begin{proof}
    Use \thmref{propLinMapExtend} to define $F_B$ as required, then by \thmref{corLinPigeon} we only need to show surjectivity.
\end{proof}
\begin{definition}{Coordinate vector}
    Given a finite vector space $V$ and an ordered basis $B$, the \underline{coordinate vector} of an element $v \in V$ with respect to $B$ is the vector in $\F^{\dim(V)}$ with its elements the coefficients of the linear combination of $v$ using $B$.
\end{definition}
\begin{corollary}
    If $|V| = |W|$ then $V \isom W$.
    \label{corSameSizeIsom}
\end{corollary}
\begin{proof}
    Consider the isomorphisms via $\F^{|V|}$ generated by \thmref{corCoordinateVectorMap}.
\end{proof}
\section{Direct Sums}
\begin{definition}{(External) direct sum}
    For vector spaces $V$ and $W$ over a field $\F$, the \underline{(external) direct sum} $V \oplus W$ is the vector space $V \times W$ with operations:
    \begin{itemize}
        \item $(v_1, w_1) + (v_2, w_2) = (v_1+v_2, w_1 + w_2)$
        \item $\lambda(v, w) = (\lambda v, \lambda w)$
    \end{itemize}
\end{definition}
\begin{lemma}
    For $V, W$ as above,
    \begin{equation*}
        \dim(V \oplus W) = \dim(V) + \dim(W)
    \end{equation*}
    \label{lemSizeDirSum}
\end{lemma}
\begin{proof}
    We can prove this by finding an isomorphism. Consider the map from the two coordinate vectors for $u$ and $w$ and mapping to one coordinate vector (one stacked on the other). Prove that this is an isomorphism.
\end{proof}
\begin{proposition}
    Let $V$ be a vector space over $\F$ and $U, W \leq V$. There exists a surjective linear map:
    \begin{align*}
        \phi : U \oplus W &\mapsto U + W\\
        (u, w) &\mapsto u + w
    \end{align*}
    with $\ker(\phi) = U \cap W$.
    \label{propDirSumMapSum}
\end{proposition}
\begin{proof}
    Linearity and surjectivity are easy. For the kernel, show that $x \mapsto (x, -x)$ is an isomorphism of $U \cap W$ and the kernel.
\end{proof}
\begin{corollary}[Sum-intersection formula]
    Let $V$ be a finite dimensional vector space over $\F$ with $U, W \leq V$:
    \begin{equation*}
        \dim(U + V) = \dim(U) + \dim(W) - \dim(U \cap W)
    \end{equation*}
    \label{corSumIntersection}
\end{corollary}
\begin{proof}
    Use \thmref{propDirSumMapSum} and \thmref{thmRankNullity}.
\end{proof}
\begin{definition}{(Internal) direct sum}
    Let $U, W \leq V$. Then $U \oplus W$ is the \underline{(internal) direct sum} of $V$ if:
    \begin{enumerate}
        \item $U + W = V$
        \item $U \cap W = \{\zv\}$
    \end{enumerate}
    and we write $V = U \oplus W$.
\end{definition}
\begin{definition}{Direct complement}
    For $U \leq V$, a \underline{direct complement} to $U$ in $V$ is a subspace $W \in V$ such that:
    \begin{equation*}
        V = U \oplus W
    \end{equation*}
\end{definition}
\begin{proposition}
    The direct complement always exists.
    \label{propDirCompExists}
\end{proposition}
\begin{proof}
    Argue by bases.
\end{proof}
\end{document}